%% ma: L12-B14-0053
%% lop: 12
%% chuong: 5
%% bai: 14
%% dang: 12.14.1
%% loai: TN4
%% muc: NB
%% dap_an: "C"
%% nguon: "To_n_12_PT_M_t_Ph_ng.pdf (Đề kiểm tra đánh giá kiến thức, Chuyên đề Phương trình mặt phẳng, Nguyễn Anh Khôi - Lê Minh Khôi), Câu 7"
%% kien_thuc: "Bài 14 (phương trình tổng quát của mặt phẳng, vectơ pháp tuyến)"
%% trang_thai: chinh-thuc
%% ghi_chu: "Tài liệu gốc không kèm đáp án; đáp án do Claude giải và kiểm bằng sympy."
\begin{ex}
Trong không gian $Oxyz$, cho mặt phẳng $(\alpha)\colon 2x-3y-4z+1=0$. Khi đó, một vectơ pháp tuyến của $(\alpha)$ là
\choice{$\vec n=(2;3;-4)$}{$\vec n=(2;-3;4)$}{\True $\vec n=(-2;3;4)$}{$\vec n=(-2;3;1)$}
\loigiai{$(\alpha)$ có vectơ pháp tuyến $(2;-3;-4)$, nên $(-2;3;4)$ cũng là vectơ pháp tuyến. Chọn C.}
\end{ex}
